\chapter{绪论}

\section{研究背景与意义}

\subsection{研究背景}

随着人工智能、传感器技术和高性能计算的飞速发展，自动驾驶技术已从实验室概念验证阶段逐步迈入大规模道路测试和商业化落地阶段。根据国际自动机工程师学会（Society of Automotive Engineers, SAE）的分级标准，自动驾驶系统分为L0至L5共六个等级\cite{paden2016}，其中L3及以上等级要求系统能够在特定或全部场景下独立完成驾驶任务。路径规划模块作为自动驾驶系统感知—决策—控制链路中的核心环节，负责根据环境信息生成安全、高效、舒适的行驶轨迹，其性能直接决定了自动驾驶系统的安全水平和用户体验。

在实际城市道路环境中，自动驾驶车辆面临的一个核心挑战是多障碍物密集场景的处理。城市交叉路口、繁忙路段、施工区域等典型场景中，往往同时存在大量的动态障碍物，如机动车、行人、非机动车，以及静态障碍物，如施工围挡、违停车辆、临时路障等。这些障碍物数量多、分布复杂、运动状态多变，对路径规划算法提出了极高的要求\cite{gonzalez2016}。具体而言，多障碍物场景下的路径规划需要解决以下几个关键问题。

第一，可行域压缩问题。当多个障碍物同时出现在自车前方时，每个障碍物都会在规划空间中占据一定区域，这些禁行区域的叠加可能导致剩余可行域被严重压缩，甚至收缩为零，使规划算法无法找到可行解。

第二，决策耦合问题。对于每个障碍物，规划系统需要做出跟随、超越、避让等决策，而这些决策之间往往存在耦合关系。例如，对障碍物A选择左绕可能导致与障碍物B发生冲突，需要全局考虑所有障碍物的决策组合才能找到最优方案。

第三，实时性问题。多障碍物场景下的规划问题复杂度随障碍物数量急剧增加，而自动驾驶系统通常要求规划模块在100毫秒内完成一个规划周期\cite{fan2018baidu}，这对算法的计算效率提出了严格的约束。

百度Apollo作为全球领先的开源自动驾驶平台，提供了从感知、预测到规划、控制的完整技术栈\cite{apollo2024}。其Planning模块采用了基于凸优化的路径规划方法，通过Scenario-Stage-Task三级管理架构实现了场景化的规划策略，具有较强的工程实践价值和学术研究价值。然而，通过对Apollo Planning模块源码的深入分析，本文发现其在多障碍物密集场景下仍存在一系列亟待解决的问题，包括PathBoundsDecider逐障碍物独立决策导致的矛盾避让方向、动态障碍物被路径决策完全排除、以及安全裕度对所有障碍物采用统一固定值而缺乏基于威胁度的分级机制等。这些问题直接影响了系统在复杂交通环境中的通行效率和规划质量。

\subsection{研究意义}

本课题立足于百度Apollo这一国内事实标准的开源自动驾驶平台，针对多障碍物密集场景下的路径规划失败问题展开研究，其意义体现在工程落地、学术推进与方法论迁移三个层面。

\textbf{工程意义}\quad{}百度Apollo自2017年开源以来，已发展为全球装机量最大的自动驾驶开源栈之一，其Planning模块几乎成为国内L4级自动驾驶研发的事实参考实现。蘑菇车联、元戎启行、轻舟智航等量产公司的规划栈或直接fork自Apollo，或在Apollo的Scenario-Stage-Task架构基础上二次开发，百度自家的萝卜快跑Robotaxi亦已在武汉、北京等多座城市开展无人化商业运营。这意味着针对PathBoundsDecider缺陷的任何切实改进，都具备直接下沉到现网车队与量产产品的可能，而不是停留在仿真实验的工具链演示。本文所有改进均以patch形式嵌入Apollo原有数据流，不引入外部求解器与额外通信开销，单个规划周期计算增量控制在毫秒级，从工程可移植性上满足下游企业的低成本集成需求。

\textbf{学术意义}\quad{}路径-速度解耦架构自DARPA Urban Challenge以来长期作为工业级规划系统的主流范式，其核心缺陷在于路径阶段的静态化假设与速度阶段的时序信息之间存在结构性断层。本文识别的矛盾避让方向、动态障碍物路径失声、统一安全裕度三类问题，本质上都是同一断层在不同算子上的投影。通过引入到达时间函数$\tau(s)$将速度维度的时序预测反哺到路径边界构建阶段，本文为路径-速度解耦框架在动态场景下的信息共享给出了一个可操作的改进范式，既保留了解耦架构在求解效率与凸性上的优势，又弥补了其在多动态障碍物场景下的决策盲区，可为后续相关架构的改进提供参考思路。

\textbf{方法论意义}\quad{}本文构建的多因子威胁度评估、基于组合优化的全局避让方向决策、时空走廊融合三者耦合的技术路线，其核心思想不局限于Apollo这一具体平台。多因子威胁度将距离、相对速度、碰撞时间等异质风险信号在可解释权重下聚合为统一标量，可直接替换其他规划系统中基于固定阈值的粗粒度危险等级；基于重叠分组的避让方向枚举将原本的NP形式决策压缩到线性规模，其最优性证明对任意以横向间隙为代价的路径规划器均成立；走廊融合所采用的时变边界表达可平移到基于优化的局部规划、基于采样的Frenet格点规划以及混合整数规划等多种技术栈。上述组合思路因此具备跨平台、跨算法范式的迁移价值，为多障碍物场景下的规划问题提供了一种超越具体实现的通用处理框架。

\section{国内外研究现状}

\subsection{路径规划方法分类}

自动驾驶路径规划技术经过数十年发展，已形成了多种成熟的算法体系，可大致分为以下四类\cite{gonzalez2016,paden2016,katrakazas2015realtime,schwarting2018planning,zhu2024decision,li2019icvmp}。四类方法的层次关系与代表算法如图\ref{fig:method_taxonomy}所示。图中以粉色高亮的优化方法分支即本文所属技术路线，Apollo的PiecewiseJerkPathOptimizer（PJPO）与本文改进算法均属于该分支。

\begin{figure}[H]
\centering
\resizebox{\textwidth}{!}{% 路径规划方法分类树
\begin{tikzpicture}[
    scale=1.0, transform shape,
    lvl1/.style={rectangle, rounded corners=3pt, draw=deepblue, fill=deepblue, text=white, minimum width=3.2cm, minimum height=0.9cm, font=\normalsize\bfseries, align=center},
    lvl2/.style={rectangle, rounded corners=2pt, draw=deepblue, fill=lightblue, minimum width=2.6cm, minimum height=0.8cm, font=\small, align=center},
    methodbox/.style={rectangle, rounded corners=2pt, draw=deepblue!60, fill=bglight, minimum width=2.6cm, font=\footnotesize, align=center, inner sep=6pt},
    methodboxhl/.style={rectangle, rounded corners=2pt, draw=deeppink, fill=lightpink!50, minimum width=2.6cm, font=\footnotesize, align=center, inner sep=6pt},
    hl/.style={rectangle, rounded corners=2pt, draw=deeppink, fill=improvedpink, minimum width=2.6cm, minimum height=0.8cm, font=\small, align=center},
    arr/.style={-Stealth, thick, deepblue},
    arrhl/.style={-Stealth, thick, deeppink},
]

\node[lvl1] (root) at (0,0) {路径规划方法};

\node[lvl2] (g1) at (-6,-1.6) {图搜索方法};
\node[lvl2] (g2) at (-2,-1.6) {采样方法};
\node[hl]   (g3) at (2,-1.6) {优化方法};
\node[lvl2] (g4) at (6,-1.6) {学习方法};

\node[methodbox] (b1) at (-6,-3.8) {A*\\Dijkstra\\D* / Hybrid A*};
\node[methodbox] (b2) at (-2,-3.8) {RRT\\RRT*\\PRM};
\node[methodboxhl] (b3) at (2,-3.8) {QP / iLQR\\Apollo PJPO\\MPC};
\node[methodbox] (b4) at (6,-3.8) {强化学习\\模仿学习\\Transformer};

\draw[arr] (root) -- (g1);
\draw[arr] (root) -- (g2);
\draw[arrhl] (root) -- (g3);
\draw[arr] (root) -- (g4);

\draw[arr] (g1) -- (b1);
\draw[arr] (g2) -- (b2);
\draw[arrhl] (g3) -- (b3);
\draw[arr] (g4) -- (b4);

\node[font=\footnotesize, deeppink] at (2,-5.5) {本文工作所属分支};

\end{tikzpicture}
}
\caption{路径规划方法分类}
\label{fig:method_taxonomy}
\end{figure}

如图所示，四类方法各有其适用边界。图搜索方法依赖状态空间离散化，适合低维、低速场景，其完备性保证使其成为泊车与OpenSpace规划的首选。采样方法通过随机采样规避维度灾难，在复杂约束下具有较强的可行性发现能力，但最优性依赖采样密度，在多障碍物密集场景下存在组合爆炸风险。优化方法将规划问题显式表达为带约束的数学规划问题，能在凸性假设下保证全局最优与实时性，是Apollo、Waymo等工业级L4平台的主流选择，但对问题的非凸性处理能力有限。学习方法近年受到广泛关注，但由于决策过程不可解释与形式化安全验证困难，目前主要作为传统方法的补充。

本文的技术栈定位明确选择优化方法分支。一方面，优化方法在Apollo平台已具备成熟工程基础，改进成果可直接复用求解器与数据结构；另一方面，本文识别的PathBoundsDecider缺陷本质上是凸化过程中丢失全局信息所致，属于优化范式内部可解决的问题，无需引入额外的采样或学习开销。

\begin{enumerate}
\item \textbf{图搜索方法}

图搜索方法将连续的规划空间离散化为图结构，通过搜索算法寻找最优路径。经典的A*算法\cite{hart1968astar}通过引入启发式函数指导搜索方向，在保证最优性的同时显著提升了搜索效率。\citet{dolgov2008hybrid,dolgov2010}提出的Hybrid A*算法将连续状态空间的车辆运动学约束与离散图搜索相结合，能够生成满足非完整约束的可行路径，被广泛应用于自动驾驶的泊车和低速场景。该算法在Apollo的OpenSpace Planner中得到了实际应用。图搜索方法的优点是理论基础扎实、完备性好，但在高维状态空间中面临计算复杂度快速增长的问题。

\item \textbf{采样规划方法}

采样规划方法通过在状态空间中随机采样来构建搜索树或路线图，避免了对状态空间的显式离散化。快速探索随机树（Rapidly-exploring Random Tree, RRT）\cite{lavalle1998rrt}及其改进版本RRT*\cite{karaman2011sampling}与双向扩展的RRT-Connect\cite{kuffner2000rrtconnect}是该类方法的代表，能够高效处理高维规划问题；\citet{kuwata2009rrt}将其拓展到城市道路实时规划，在2007年DARPA Urban Challenge的MIT车队与CMU的Boss车队\cite{urmson2008boss}中得到实战验证。\citet{werling2010}提出了基于Frenet坐标系的轨迹规划方法，通过在参考线的横向和纵向空间中独立采样终端状态，使用多项式曲线连接起始状态和终端状态，生成候选轨迹集合并根据代价函数进行评估筛选；\citet{li2022frenet}进一步从运输系统视角论证了Frenet框架在结构化道路下的普适性；\citet{mcnaughton2011lattice}则给出了面向高速公路的时空状态格（Spatiotemporal State Lattice）扩展。这一方法论被Apollo的Lattice Planner所采用。采样方法具有灵活性强、可处理复杂约束的优点，但在多障碍物密集场景下面临采样效率低和组合爆炸的问题。

\item \textbf{优化规划方法}

基于优化的方法将路径规划建模为数学优化问题，通过求解带约束的目标函数获得最优轨迹。二次规划（QP）方法将轨迹表示为分段多项式，以平滑度、偏离参考线程度等为目标函数，以障碍物约束、动力学约束等为约束条件，利用OSQP\cite{stellato2020osqp}等高效求解器求解，非线性场景则可借助IPOPT\cite{wachter2006ipopt}等内点法求解器。该路线的早期代表是\citet{borrelli2005mpc}提出的基于模型预测控制（Model Predictive Control, MPC）的主动转向方法、\citet{xu2012realtime}提出的面向自动驾驶的实时轨迹优化器，以及\citet{ziegler2014}在Bertha Benz自动驾驶车上首次将连续局部优化部署到公开城市道路的工作。\citet{fan2018baidu}详细描述了Apollo期望最大化（Expectation-Maximization, EM）Motion Planner的设计，采用路径QP和速度QP的两阶段优化方法，\citet{zhou2021pathqp,zhou2021dliaps}进一步给出了Apollo路径QP的分段Jerk形式化表达与双环迭代锚点平滑（DL-IAPS）和分段Jerk速度优化（PJSO）算法。\citet{chen2019em}提出了迭代EM规划框架，进一步增强了优化方法的灵活性；\citet{sun2022robust}则专门针对大曲率道路给出了鲁棒路径优化器。受无人机轨迹生成中最小Snap方法\cite{mellinger2011minsnap}与安全飞行走廊（Safe Flight Corridor）\cite{liu2017sfc}的启发，地面自动驾驶也逐步引入了类似的走廊凸分解思想。优化方法的优点是能够生成高质量的平滑轨迹，并在凸优化框架下保证全局最优性，但对问题的凸性要求较高，多障碍物场景下的非凸约束处理是主要挑战。

\item \textbf{学习型方法}

近年来，基于深度学习和强化学习的规划方法受到了广泛关注\cite{zhang2021autonomous}。端到端学习方法直接从传感器输入映射到控制输出，具有较强的场景泛化能力；\citet{alahi2016social}提出的Social LSTM和基于Transformer的预测性运动规划\cite{zhang2024trans}也为规划层提供了更可靠的先验；\citet{li2024decision}系统梳理了多模态协同与端到端训练的技术演进，进一步展示了该路线的发展趋势。然而，由于决策过程的不可解释性和安全性验证困难，学习型方法在工业级自动驾驶系统中的应用仍然有限，通常作为传统方法的辅助或补充\cite{liu2024e2e}；\citet{shalevshwartz2017rss}提出的责任敏感安全（Responsibility-Sensitive Safety, RSS）模型从形式化角度为学习型决策提供了安全外壳。
\end{enumerate}

\begin{table}[H]
\centering
\caption{四类路径规划方法对比}
\label{tab:method_compare}
\small
\begin{tabular}{@{}p{2.3cm}p{2.2cm}p{2.2cm}p{2.2cm}p{2.3cm}@{}}
\toprule
\textbf{维度} & \textbf{图搜索} & \textbf{采样} & \textbf{优化} & \textbf{学习} \\
\midrule
完备性 & 分辨率完备 & 概率完备 & 局部最优 & 无形式保证 \\
最优性 & 离散最优 & 渐近最优 & 凸下全局最优 & 数据依赖 \\
实时性 & 随维度恶化 & 高 & 高（10--100\,ms） & 推理快、训练贵 \\
非凸处理 & 强 & 强 & 弱（需凸化） & 自然支持 \\
工程成熟度 & 已工业化 & 已工业化 & L4主流 & 仍在研 \\
Apollo对应 & OpenSpace & Lattice Planner & PublicRoad (PJPO) & 辅助模块 \\
\bottomrule
\end{tabular}
\end{table}

表\ref{tab:method_compare}揭示了本文技术选型的底层逻辑。优化方法在实时性与工程成熟度上取得了最好的平衡，是Apollo PublicRoad Planner的核心范式；但其非凸处理弱的短板恰恰在多障碍物密集场景下被放大，PathBoundsDecider的逐障碍物贪心决策正是典型表现。本文的改进抓手落在优化范式内部，通过时变最优通行带的组合优化将凸化前的全局信息注回，而无需跳出优化框架。

\subsection{Apollo Planning模块技术演化脉络}

百度Apollo自2017年开源以来，其Planning模块经历了从EM Planner、PublicRoad Planner到OpenSpace Planner的三阶段演化，理解这一脉络是识别PathBoundsDecider定位的前提。

\textbf{EM Planner时期，2017至2019年}\quad{}该阶段采用期望最大化思想在Frenet坐标系下交替求解路径与速度子问题\cite{fan2018baidu}。路径阶段以动态规划搜索粗解、以QP平滑精解，速度阶段在ST图上重复同一套DP粗解加QP精解的模式。EM Planner首次将路径-速度解耦的两阶段优化完整落地到工业级L4栈，成为后续所有Apollo Planner的范式基础。

\textbf{PublicRoad Planner时期，2019年至今}\quad{}该阶段引入Scenario-Stage-Task三级管理架构，对LaneFollow、Lane Change、Park and Go等场景进行显式建模。路径层由\citet{zhou2021pathqp}的分段Jerk路径优化器（PiecewiseJerkPathOptimizer, PJPO）接管，速度层对应为PJSO\cite{zhou2021dliaps}。其中PathBoundsDecider作为PJPO的前置Task，负责将障碍物决策转化为路径QP的左右边界约束，直接决定了优化空间的形状，也是本文改进的核心切入点。

\textbf{OpenSpace Planner补充时期，2020年至今}\quad{}为支持泊车、OpenSpace等结构化先验失效的场景，Apollo引入Hybrid A\*\cite{dolgov2008hybrid}与DL-IAPS\cite{zhou2021dliaps}的自由空间规划链路，与PublicRoad Planner并行。

基于该演化脉络可以明确本文定位。本文不触碰EM与PJPO的优化器本身，而是在PathBoundsDecider这一路径边界构建环节把全局决策能力注入进去，属于在PublicRoad主干内部做精细化升级，改动颗粒度小，可与上下游解耦回归。

\subsection{Apollo平台相关研究}

除上述平台自研工作外，Apollo还吸引了大量的外部学术研究和工程实践。

在算法分析方面，\citet{kessler2023}在Apollo规划栈中实现了混合整数运动规划，用于处理德国道路上的复杂驾驶场景。\citet{li2023quant}对Apollo规划参数化对关键场景中轨迹规划的影响进行了定量分析，揭示了参数选择对规划质量的显著影响。\citet{li2023valid}基于Apollo平台验证了自动驾驶功能，展示了仿真与实车测试的结合方法。

在国内研究方面，\citet{gao2024apollo}探讨了基于Apollo平台的自动驾驶教学创新方法。\citet{sun2023urban}研究了面向城市工况的运动规划方法。\citet{ma2023parking}将智能优化算法应用于Apollo框架下的泊车路径规划。在仿真验证方面，CARLA\cite{dosovitskiy2017carla}、LGSVL\cite{rong2020lgsvl}等开放仿真器以及nuScenes\cite{caesar2020nuscenes}等大规模数据集为Apollo类平台提供了可复现的评估基座。这些研究为Apollo平台的实际应用提供了丰富的参考。

\subsection{静态多障碍物场景规划研究}

当多个静态障碍物同时出现在自车前方时，其核心挑战是将多约束非凸可行域转化为可解的数学结构。

\textbf{走廊凸分解}\quad{}\citet{wang2023corridor}提出了基于时空走廊的非结构化环境轨迹规划方法，通过构建安全通道来简化多障碍物约束处理；\citet{yang2025stcorridor}在三维时空走廊内以分段贝塞尔曲线求解最优轨迹，在安全性与灵活性上相较时间弹性带算法取得了明显提升；\citet{qiao2025stjoint}则在结构化道路下引入时间-纵向-横向三维时空决策并与时空安全走廊的二次规划相结合；\citet{xu2025interaction}在城市交互场景中将部分可观测马尔可夫决策过程（Partially Observable Markov Decision Process, POMDP）行为决策与基于约束迭代线性二次调节器（Constrained Iterative Linear Quadratic Regulator, CILQR）的时空安全走廊求解进行耦合。上述工作的共同点是将多障碍物约束统一抽象为走廊边界，进而将非凸问题转化为凸问题。

\textbf{多障碍物交互建模}\quad{}\citet{chen2024avoid}研究了复杂道路环境下智能车辆的横纵向避障问题，提出了考虑障碍物交互关系的综合避障策略；\citet{guo2024survey}对自动驾驶场景和场景数据库进行了全面综述，为多障碍物场景的构建和评估提供了参考框架。

然而上述走廊方法普遍预设走廊已经给定，并未回答多障碍物之间应如何选择左绕或右绕从而最大化走廊宽度这一前置问题。Apollo的PathBoundsDecider正是因为在这一前置步骤采用逐障碍物独立决策的贪心策略，才导致后续走廊被压缩甚至阻塞，这构成本文的直接立题。

\subsection{动态障碍物的时间维度处理}

相比静态场景，动态障碍物引入了时间维度，其处理方式演进呈现出清晰的脉络。

\textbf{路径-速度解耦阶段}\quad{}\citet{kant1986pvd}于1986年提出的路径-速度解耦（Path-Velocity Decomposition, PVD）是整个领域的理论起点，先在笛卡尔空间规划路径绕开静态障碍物，再在路径-时间（Path-Time, PT）空间上规划速度规避动态障碍物。这一范式的优点是将高维时空规划降维到两个二维子问题，在Apollo EM Planner\cite{fan2018baidu}和PublicRoad Planner中均有体现，ST图本质上就是Kant-Zucker PT空间的直接继承。

\textbf{ST图与局部时间优化阶段}\quad{}在PVD基础上，ST图被广泛用于速度规划。\citet{liu2017speed}进一步提出基于时间优化的速度剖面规划方法，将时间维度显式纳入优化变量；\citet{vandenberg2008rvo}提出的互惠速度障碍（Reciprocal Velocity Obstacle, RVO）为多智能体动态避让提供了速度层的几何刻画；\citet{ziegler2014}在Bertha项目中将动态障碍物作为时变约束集成进局部连续优化器。

\textbf{时间弹性带与时空联合阶段}\quad{}\citet{rosmann2012teb,rosmann2013sparse}提出的时间弹性带（Timed Elastic Band, TEB）将时间显式作为优化变量，在局部轨迹层面整合时间维度，是PVD的第一次重大突破。近年的研究进一步从速度规划层向上反哺路径规划，\citet{yang2025stcorridor}、\citet{qiao2025stjoint}的时空走廊工作已隐含路径-速度耦合的思想，\citet{li2025perturbation}则在此基础上引入了障碍物扰动鲁棒性。

\textbf{本文在脉络中的定位}\quad{}尽管上述工作推进了动态障碍物处理的时间维度，但无一例外地仍将路径决策阶段的动态障碍物视为当前时刻的静态快照，Apollo PathBoundsDecider即是该假设的典型产物。本文的贡献在于把到达时间函数引入路径边界构建阶段，将动态障碍物未来时刻将离开当前位置的预测信息直接注入路径决策，实现路径层对时间维度的显式建模，而非仅依赖速度层事后补偿。

\subsection{现有研究不足}

综合分析现有研究，可以发现以下不足。

第一，针对Apollo Planning模块路径边界决策在多障碍物场景下的缺陷缺乏系统性分析。现有研究多关注速度规划层面的参数调优或单一模块的改进，而对路径边界构建阶段（PathBoundsDecider）逐障碍物独立决策导致的矛盾避让方向问题缺乏深入研究。

第二，多障碍物的避让方向选择缺乏全局最优性保证。Apollo对每个障碍物独立选择左绕或右绕的贪心策略，在多障碍物场景下可能导致相邻障碍物被分配矛盾的避让方向，进而导致路径不可行。该问题的本质是一个组合优化问题，现有研究尚未从数学优化的角度给出系统性的解决方案。

第三，路径决策未考虑动态障碍物的时间演化。Apollo的路径边界构建仅处理静态障碍物，动态障碍物被完全交由速度规划阶段处理。这种割裂使得路径决策无法利用“动态障碍物在未来时刻将离开当前位置”的预测信息来扩展可行路径空间。此外，\citet{li2025perturbation}指出现有时空联合规划算法普遍假设障碍物位置精确已知，对感知扰动缺乏鲁棒应对机制，这一结论同样适用于Apollo当前的路径边界决策。

本文将针对上述不足，提出系统性的改进方案。

\section{研究内容与技术路线}

\subsection{研究内容}

本文的研究内容主要包括以下五个方面。

\begin{enumerate}
\item \textbf{Apollo Planning模块深度分析}\quad{}对Apollo 11.0版本的Planning模块进行源码级深入分析，理解其Scenario-Stage-Task三级管理架构、ReferenceLine平滑机制、PublicRoadPlanner中源自早期EM Planner思想的路径-速度解耦两阶段优化流程、PathBoundsDecider路径边界构建和SpeedBoundsDecider的ST边界构建机制。

\item \textbf{多障碍物场景问题识别}\quad{}通过源码分析和数学建模，系统识别Apollo现有算法在多障碍物场景下的三个核心问题，即安全裕度一刀切压缩通行空间、PathBoundsDecider逐障碍物独立决策导致矛盾避让方向、以及动态障碍物被路径决策完全排除导致路径-速度协调缺失。三个问题分别从安全距离、决策粒度和时空协调三个维度指向改进方向。

\item \textbf{时变最优通行带问题建模}\quad{}将多障碍物场景下的避让方向选择建模为组合优化问题，引入到达时间函数将动态障碍物的预测轨迹纳入路径决策，建立统一的时变最优通行带数学框架。将时变通行带投影至ST空间解读为时空可通行走廊，建立路径决策与速度规划之间的约束传递机制。

\item \textbf{基于重叠分组的最优求解算法设计与实现}\quad{}通过引理证明最优分割的连续性、通过定理证明最大间隙策略的最优性，提出基于扫描线的$O(n \log n)$时变障碍物分组算法，在Apollo平台上实现并集成到PathBoundsDecider中。

\item \textbf{仿真验证与对比分析}\quad{}在Apollo仿真环境中设计多组多障碍物测试场景，重点验证”Baseline路径阻塞而改进方案可通行”的核心对比，从规划成功率、通行效率、停车次数等维度进行定量分析。
\end{enumerate}

\subsection{技术路线}

本文的技术路线如图\ref{fig:tech_route}所示，分为七个阶段依次推进。

\begin{figure}[H]
\centering
% 论文版技术路线图（七步研究路线，对齐 Ch4/Ch5/Ch6 双层框架）
\begin{tikzpicture}[
    node distance=0.35cm,
    scale=1.05, transform shape,
    block/.style={rectangle, draw=deepblue, fill=skyblue, text width=10cm, text centered, rounded corners, minimum height=0.9cm, font=\normalsize},
    coreblock/.style={rectangle, draw=deeppink, fill=improvedpink, text width=10cm, text centered, rounded corners, minimum height=0.9cm, font=\normalsize},
    endblock/.style={rectangle, draw=deepblue, fill=outputpink, text width=10cm, text centered, rounded corners, minimum height=0.9cm, font=\normalsize},
    arrow/.style={thick, -Stealth, deepblue}
]

\node (s1) [block] {\textbf{Apollo Planning 源码级分析}};

\node (s2) [block, below=of s1] {\textbf{多障碍物场景局限识别}};

\node (s3) [block, below=of s2] {\textbf{多障碍物统一规划统一建模}};

\node (s4) [coreblock, below=of s3] {\textbf{最优通行带求解算法设计}};

\node (s5) [coreblock, below=of s4] {\textbf{时空可通行走廊与速度协调}};

\node (s6) [block, below=of s5] {\textbf{Apollo 集成实现}};

\node (s7) [endblock, below=of s6] {\textbf{多场景仿真对比与验证}};

\draw[arrow] (s1) -- (s2);
\draw[arrow] (s2) -- (s3);
\draw[arrow] (s3) -- (s4);
\draw[arrow] (s4) -- (s5);
\draw[arrow] (s5) -- (s6);
\draw[arrow] (s6) -- (s7);

\end{tikzpicture}

\caption{技术路线图}
\label{fig:tech_route}
\end{figure}

\section{论文组织结构}

本文共分为七章，各章内容安排如下。

第一章为绪论，介绍研究背景与意义、国内外研究现状、研究内容与技术路线。

第二章为相关技术与理论基础，介绍Frenet坐标系、凸优化与二次规划、SL图与路径边界构建、Apollo平台架构等基础知识。

第三章为Apollo Planning模块深度分析，对Planning模块的总体架构、ReferenceLine平滑、路径-速度解耦两阶段优化流程以及PathBoundsDecider路径边界决策逻辑进行源码级分析。

第四章为多障碍物场景下的问题分析，从安全距离、决策粒度和时空协调三个维度系统剖析Apollo PathBoundsDecider的规划缺陷，即安全裕度一刀切压缩通行空间、逐障碍物独立决策导致矛盾避让方向、以及动态障碍物被完全排除导致路径-速度协调缺失。

第五章为改进算法设计与实现，建立时变最优通行带的数学模型，从SL层提出多因子威胁度评估与差异化安全距离、基于重叠分组的最大间隙策略并证明其最优性，从ST层引入到达时间函数和预测轨迹将动态障碍物纳入路径决策并构建时空可通行走廊与速度规划协调机制，最后给出与Apollo系统的集成方案。

第六章为实验与结果分析，介绍实验环境、场景设计、评价指标，展示对比实验结果并进行分析。

第七章为总结与展望，总结全文工作和贡献，讨论研究局限性和未来方向。
